\documentclass[11pt]{article}
\usepackage[a4paper,margin=1in]{geometry}
\usepackage{amsmath,amssymb,amsthm}
\usepackage{booktabs}
\usepackage[hidelinks]{hyperref}

\newtheorem{theorem}{Theorem}
\newtheorem{lemma}[theorem]{Lemma}
\newtheorem{corollary}[theorem]{Corollary}
\theoremstyle{definition}
\newtheorem{definition}[theorem]{Definition}
\newtheorem{observation}[theorem]{Observation}
\theoremstyle{remark}
\newtheorem{remark}[theorem]{Remark}

\newcommand{\R}{\mathbb{R}}
\newcommand{\N}{\mathbb{N}}
\newcommand{\F}{\mathcal{F}}
\newcommand{\lean}[1]{\texttt{#1}}

\title{Twin primes as a positivity statement:\\
a machine-checked reformulation and three computed forms\\ of the parity barrier}
\author{Seth Caldwell\thanks{Repository: \url{https://github.com/oceanseth/primefold}.
The Lean development, the experiments and the text were produced in open sessions with
the language model Claude Fable~5.2 (Anthropic); the full turn history is in the
repository. See the statement on AI use at the end.}}
\date{21 November 2026}

\begin{document}
\maketitle

\begin{abstract}
We do not prove the twin prime conjecture. We record what a short, fully logged
attempt produced, positive and negative. (1)~A Lean~4/Mathlib proof, with no
hypotheses, that the twin prime conjecture is equivalent to the strict
positivity, for arbitrarily large $x$, of an explicit weighted integral of
$|P_x|^2$, where $P_x$ is a Dirichlet polynomial over the primes in
$(x, \tfrac{51}{50}x)$. The underlying identity holds for every finite set of
integers, which turns it into a test: an argument for positivity must fail on
sets that have no two elements differing by~$2$. (2)~Three computations with such
decoy sets near $10^8$. The Hardy--Littlewood local model predicts the number of
twin primes to $1.3\%$ and predicts thousands of pairs for sets that contain none.
Splitting the primes $p$ by the parity of $\Omega(p+2)$ gives two halves that are
indistinguishable in size and in residue classes to every prime modulus below
$500$, one with no pairs and one with twice the predicted number. Bilinear sums
over factorisations of $p+2$ separate the halves only when both coefficient
sequences are the Liouville function. (3)~A formal statement of the conditional
chain from a uniform zeta-ratios hypothesis to twin primes, with the final
implication and the positivity of the twin prime constant machine-checked. The
conclusion is that each route we tried ends at the same missing input, a
correlation between $\Lambda(n)$ and $\lambda(n+2)$, and that a claimed proof of
the Riemann Hypothesis, even if correct, does not supply it.
\end{abstract}

\section{Introduction}

The twin prime conjecture asserts that there are infinitely many primes $p$ with
$p+2$ prime. This note began with a narrower question. On 20 November 2026 a
Lean archive was published claiming a formal proof of the Riemann Hypothesis and
of its analogue for Dirichlet $L$-functions \cite{rhclaim}. At the time of
writing the archive has been compiled and re-checked independently by
A.~Sutherland \cite{rhcheck}, with some kernel checks still running; the
mathematics has had no expert reading. We have not built or checked the archive
ourselves; we have only searched its source text (Section~\ref{sec:discussion}).
The question put to us was whether such a proof, if it stands, could
be extended to the twin prime conjecture.

The answer is no, for a reason that is standard in sieve theory and that the
computations below make concrete. Information about the location of zeros
controls primes in arithmetic progressions one modulus at a time. On average over
moduli that level of control is already unconditional (Bombieri--Vinogradov), and
even the stronger Elliott--Halberstam conjecture gives only bounded gaps: at most
$12$ \cite{maynard}, or $6$ in its generalised form \cite{polymath}. The
obstruction to reaching $2$ is Selberg's parity phenomenon \cite{selberg}: sieve
methods fed with distributional information cannot distinguish integers with an
even number of prime factors from those with an odd number.

Rather than stop at that sentence, we did three things, all of which are in the
repository and all of which can be re-run.

\begin{enumerate}
\item We formalised an exact reformulation of the conjecture as a positivity
statement (Section~\ref{sec:lean}). It is a restatement and has the same
difficulty as the conjecture; its use is as a filter for proposed arguments.
\item We built that filter and ran it (Sections~\ref{sec:decoy}--\ref{sec:type2}).
\item We formalised the statement of a recent conditional result of Lee
\cite{lee}, which derives the Hardy--Littlewood asymptotic for prime pairs from
a uniform version of the Conrey--Snaith ratios conjecture \cite{conreysnaith},
and machine-checked the parts of the chain that are elementary
(Section~\ref{sec:chain}).
\end{enumerate}

Nothing here is evidence for or against the conjecture beyond the numerical
agreement with Hardy--Littlewood \cite{hl} that has been known for a century.
The negative results are the content.

\section{The formalised reformulation}\label{sec:lean}

All statements in this section are proved in Lean~4 with Mathlib
\cite{mathlib}, in the directory \lean{lean/} of the repository. The proofs
contain no \lean{sorry}, no added axioms and no \lean{native\_decide}; the
command \lean{\#print axioms} reports only \lean{propext},
\lean{Classical.choice} and \lean{Quot.sound}. Lean names are given in
typewriter face.

We use the Fourier transform $\F\psi(u)=\int_\R \psi(v)e^{-2\pi i uv}\,dv$.

\begin{theorem}[\lean{integral\_fourier\_mul\_normSq}]\label{thm:identity}
Let $\psi:\R\to\mathbb{C}$ be continuous and integrable with $\F\psi$ integrable.
For every finite set $S\subset\N$ and real sequences $(a_n)$, $(\lambda_n)$,
\[
\int_\R \F\psi(u)\,\Bigl|\sum_{n\in S} a_n e^{-iu\lambda_n}\Bigr|^2 du
 \;=\; \sum_{m\in S}\sum_{n\in S} a_m a_n\,\psi\!\Bigl(\frac{\lambda_n-\lambda_m}{2\pi}\Bigr).
\]
\end{theorem}

This is Fourier inversion applied term by term to the expanded square. With
$\lambda_n = 2\pi x\log n$ the right side becomes
$\sum_{m,n} a_m a_n\,\psi\bigl(x(\log n-\log m)\bigr)$.

\begin{lemma}[Window lemma, \lean{eq\_add\_of\_window}]\label{lem:window}
Let $h\ge 1$ be an integer, $x>0$ real, and $m<n$ integers with
$x<m$ and $n<\bigl(1+\tfrac{1}{25h}\bigr)x$. If
$\bigl|x(\log n-\log m)-h\bigr|<\tfrac14$, then $n=m+h$.
\end{lemma}

The proof uses the elementary bounds $1-1/t\le\log t\le t-1$, so no ``for
sufficiently large $x$'' is needed.

\begin{definition}[\lean{PairTest}]
A \emph{test function for difference $h$} is a continuous, even, nonnegative
$\rho:\R\to\R$, integrable with integrable Fourier transform, such that
$\rho(v)\neq0$ with $v\ge0$ implies $|v-h|<\tfrac14$, and $\rho(v)>0$ for
$h-\tfrac1{25}<v<h$.
\end{definition}

Such functions exist for every $h\ge1$ (\lean{exists\_pairTest}): take a smooth
bump centred at $h$ plus its reflection.

\begin{theorem}[Sign criterion, \lean{pairSum\_pos\_iff}]\label{thm:sign}
Let $\rho$ be a test function for difference $h\ge1$, let $x>0$, let
$S\subset\N$ be a finite set contained in the open interval
$\bigl(x,(1+\tfrac1{25h})x\bigr)$ and let $a_n>0$ for $n\in S$. Then
\[
\sum_{m,n\in S} a_m a_n\,\rho\bigl(x(\log n-\log m)\bigr)>0
\quad\Longleftrightarrow\quad
\text{there is } m\in S \text{ with } m+h\in S .
\]
\end{theorem}

Now specialise to $h=2$, to $S$ the set of primes in $(x,\tfrac{51}{50}x)$, and to
$a_p=\log p/\sqrt p$. Put
\[
P_x(u)=\sum_{x<p<\frac{51}{50}x}\frac{\log p}{\sqrt p}\;p^{-2\pi i x u}.
\]

\begin{theorem}[\lean{exists\_twinPrime\_iff\_integral\_pos}]\label{thm:main}
There is a test function $\rho$ for difference $2$ such that the following are
equivalent:
\begin{enumerate}
\item[(i)] there are infinitely many twin primes;
\item[(ii)] for every $N$ there is a real $x>N$ with
$\displaystyle\int_\R \F\rho(u)\,|P_x(u)|^2\,du>0$.
\end{enumerate}
The equivalence holds for every test function for difference $2$
(\lean{twinPrime\_iff\_integral\_pos}).
\end{theorem}

\begin{remark}
Theorem~\ref{thm:main} is the mechanism of Lemma~3.1 of \cite{lee}, simplified:
a sharp cut-off to an interval instead of a smooth weight, one test function
supported near $\pm2$ instead of two nonnegative weights, and an exact statement
for every $x$. With $\rho(v)=b_0(v-2)+b_0(v+2)$ for an even bump $b_0$ centred at $0$,
$\F\rho(u)=2\,\F b_0(u)\cos(4\pi u)$, so (ii)
says that $|P_x|^2$ is on average larger near half-integer values of $u$ than
between them.
\end{remark}

\begin{remark}[What the theorem is worth]
The integral in (ii) equals a weighted count of twin primes in
$(x,\tfrac{51}{50}x)$. Proving it positive is therefore exactly as hard as
finding the twin primes, and the theorem is not progress toward a proof. Its
use is the following. Theorems~\ref{thm:identity} and~\ref{thm:sign} hold for
\emph{every} finite set $S$ in the window. So if an argument deduces positivity
from some list of properties of the primes, and there is a set $S$ with the same
properties and no two elements differing by $2$, the argument is wrong. This is
the parity barrier stated as a test that can be executed.
\end{remark}

\section{The decoy filter}\label{sec:decoy}

The script \lean{experiments/decoy\_filter.py} takes the $108{,}524$ primes in
$(10^8,\,1.02\cdot10^8)$, builds comparison sets, and evaluates on each a
candidate ``argument'', that is, a rule that predicts the number of pairs at
distance $2$ from stated information. It also checks Theorem~\ref{thm:identity}
numerically at small $x$, where the integral and the pair sum agree to five
significant figures.

The rule tested is the Hardy--Littlewood local model \cite{hl}: the predicted
number of pairs for a set $S$ is computed from $|S|$ and from the distribution
of $S$ over residue classes modulo every prime $q\le47$. For comparison we also
give the prediction from $|S|$ alone, within the integers having no prime factor
$\le 47$.

\begin{table}[h]
\centering
\begin{tabular}{lrrrr}
\toprule
Set & Size & Pairs & Density only & Density + residues \\
\midrule
Primes & 108{,}524 & 7{,}704 & 5{,}889 & 7{,}805 \\
Random set (control) & 108{,}524 & 7{,}822 & 5{,}889 & 7{,}806 \\
Primes $\equiv1 \pmod 3$ & 54{,}152 & 0 & 1{,}466 & 0 \\
Primes, one of each twin removed & 100{,}820 & 0 & 5{,}082 & 6{,}292 \\
Random set, made pair-free & 108{,}524 & 0 & 5{,}889 & 7{,}227 \\
\bottomrule
\end{tabular}
\caption{Pairs at distance $2$ in $(10^8,1.02\cdot10^8)$, and two predictions.
The random sets are drawn from the integers with no prime factor $\le47$.}
\label{tab:decoy}
\end{table}

\begin{observation}
The local model is accurate on the primes ($7{,}805$ predicted, $7{,}704$
found) and on a random set with the same local statistics. It predicts
$6{,}292$ and $7{,}227$ pairs for two sets that have none. No argument that uses
only the size of the set and its distribution in residue classes to these
moduli can give a positive lower bound for twin primes.
\end{observation}

These first decoys are imperfect: deleting one member of each pair leaves an
excess in the residue classes $\pm2$ of $13$ to $15$ standard deviations when the
classes are combined over all moduli, where the primes show $-0.1$. It is far too
small to change the prediction, but it raises a precise question: can a
pair-free set always hide that trace?

\section{The parity decoy}\label{sec:parity}

It can, with no tuning. Let $\Omega(n)$ be the number of prime factors of $n$
counted with multiplicity, and split the primes $p\in(10^8,1.02\cdot10^8)$ into
\[
E=\{p:\Omega(p+2)\text{ even}\},\qquad O=\{p:\Omega(p+2)\text{ odd}\}.
\]
The set $E$ contains no two elements differing by $2$: if $p$ and $p+2$ are both
in $E$ then $p+2$ is prime, so $\Omega(p+2)=1$ and $p\notin E$.

\begin{table}[h]
\centering
\begin{tabular}{lrr}
\toprule
 & $E$ ($\Omega(p+2)$ even) & $O$ ($\Omega(p+2)$ odd) \\
\midrule
Size & 53{,}988 & 54{,}536 \\
Worst residue class, prime moduli $\le47$ & $2.67\,\sigma$ & $2.57\,\sigma$ \\
Worst residue class, prime moduli $53$ to $499$ & $3.57\,\sigma$ & $3.55\,\sigma$ \\
Combined excess in the classes $\pm2$ & $+0.05\,\sigma$ & $-0.20\,\sigma$ \\
Pairs at distance $2$ & \textbf{0} & \textbf{3{,}873} \\
Local model predicts & 1{,}929 & 1{,}973 \\
\bottomrule
\end{tabular}
\caption{The two parity halves of the primes in $(10^8,1.02\cdot10^8)$. For
comparison, the full set of primes has worst classes $1.69\,\sigma$ and
$2.76\,\sigma$, and a random set $2.55\,\sigma$ and $3.51\,\sigma$.}
\label{tab:parity}
\end{table}

\begin{observation}\label{obs:parity}
$E$ and $O$ have the same size to within $1.7$ standard deviations and the same
residue statistics as a random set for every prime modulus below $500$. One has
no pairs; the other has twice the number the local model predicts. The trace
left by the cruder decoys of Section~\ref{sec:decoy} was an artefact of their
construction.
\end{observation}

The consequence is a requirement on any proof. It must use information that
separates $E$ from $O$. The defining information is the Liouville function
$\lambda(n)=(-1)^{\Omega(n)}$ at $p+2$, and the two halves have asymptotically
equal size exactly when
\begin{equation}\label{eq:chowla}
\sum_{n\le x}\Lambda(n)\,\lambda(n+2)=o(x).
\end{equation}
Statement~\eqref{eq:chowla} is a hybrid of the Hardy--Littlewood and Chowla
conjectures \cite{chowla}. To our knowledge it is open; the known results on
correlations of $\lambda$ are for averaged or logarithmically averaged sums
\cite{mr,tao}. The observed split $53{,}988/54{,}536$ is consistent with
\eqref{eq:chowla}. Among the statistics tested here, the values of
$\lambda(p+2)$ are all that separates a set with no twin primes from one with
twice the expected number, so a proof has to control them in some form.

\section{Bilinear statistics}\label{sec:type2}

The known way past parity is bilinear (``type~II'') information, as in the
theorem of Friedlander and Iwaniec that $a^2+b^4$ is prime infinitely often
\cite{fi}. The script \lean{experiments/type2\_filter.py} asks whether such
statistics separate $E$ from $O$. For coefficient sequences $\alpha,\beta$ and a
dyadic range $m\in(M,2M]$ it computes
\[
T(S)=\sum_{p\in S}\;\sum_{\substack{mn=p+2\\ M<m\le 2M}}\alpha(m)\beta(n)
\]
and reports $T(E)-T(O)$ in standard deviations of a random split of the primes
into two sets of these sizes.

\begin{table}[h]
\centering
\begin{tabular}{lrrr}
\toprule
$\alpha(m),\ \beta(n)$ & $M=100$ & $M=1000$ & $M=10000$ \\
\midrule
$1,\ 1$ (divisor count) & $-1.3$ & $-0.3$ & $-0.6$ \\
$\lambda(m),\ 1$ & $+1.2$ & $+0.2$ & $-2.9$ \\
$1,\ \lambda(n)$ & $-31.6$ & $+1.6$ & $-0.6$ \\
random signs, both (rms of 20 draws) & $0.9$ & $0.7$ & $1.2$ \\
$\lambda(m),\ \lambda(n)$ & $+162.8$ & $+146.6$ & $+142.8$ \\
\bottomrule
\end{tabular}
\caption{Separation of the parity halves by bilinear sums over $p+2=mn$,
in standard deviations.}
\label{tab:type2}
\end{table}

\begin{observation}
Only $\alpha=\beta=\lambda$ separates the halves, and it does so trivially,
because $\lambda(m)\lambda(n)=\lambda(p+2)$ is the quantity that defines them.
The entry $-31.6$ for $(1,\lambda)$ at $M=100$ is the same information: the mean
of $\lambda(m)$ over $(100,200]$ is not zero. Generic bilinear weights see
nothing. Over all the primes in the range, $\sum_p\lambda(p+2)\tau_M(p+2)$ is
$-392$, $-111$ and $-192$ out of about $49{,}000$ terms for the three values of
$M$, consistent with the cancellation a proof would need to establish.
\end{observation}

So ``use type~II information'' is not by itself an idea for twin primes: the
bilinear sum that matters is a parity sum over $p+2$, which is
\eqref{eq:chowla} again. Friedlander and Iwaniec could bound their bilinear sums
because $a^2+b^4$ lies in the Gaussian integers, where the sequence factorises.
For twins of such primes that structure is unavailable, and this is a theorem
in the repository: if $p\ne2$ is a prime of the form $a^2+b^4$ then
$p\equiv1\pmod4$, so $p+2\equiv3\pmod4$ is not a sum of two squares
(\lean{IsFIPrime.add\_two\_not\_sq\_add\_sq}). The structure that defeats parity
for $p$ cannot also hold for $p+2$.

\section{The conditional chain from zeta}\label{sec:chain}

Lee \cite{lee} derives, for every even $h>0$, the Hardy--Littlewood asymptotic
\[
\pi_h(x):=\#\{p\le x: p+h \text{ prime}\}\sim\mathfrak{S}(h)\frac{x}{(\log x)^2}
\]
from a uniform two-over-two ratios formula for $\zeta$ with a power-saving
error (Hypothesis~2.1 of \cite{lee}; a restricted form of the conjecture of
Conrey and Snaith \cite{conreysnaith}). The paper is a preprint and states that
a language model was used in generating proofs. One of us read it in full and
recomputed its identities; we found no error, but that reading is not a referee
report, and the facts it imports from the literature were not re-derived. The
record is in \lean{lean/LEE-CHECK.md}.

In Lean we state the hypothesis (\lean{RatiosHypothesis}), state Lee's theorem
as a proposition that is neither proved nor assumed as an axiom
(\lean{LeeTheorem}), and prove the following.

\begin{theorem}[\lean{infinite\_pairs\_of\_hardyLittlewood}]
If $\pi_h(x)(\log x)^2/x\to S$ with $S>0$, then for every $N$ there is a prime
$p>N$ with $p+h$ prime.
\end{theorem}

\begin{theorem}[\lean{twinConstant\_pos}]\label{thm:c2}
The twin prime constant $C_2=\prod_{p>2}\bigl(1-(p-1)^{-2}\bigr)$ is positive.
\end{theorem}

\begin{proof}[Proof sketch]
Each factor is positive, since $(p-1)^2\ge4$ for $p\ge3$. The deviations from
$1$ are bounded by $4/p^2$ and hence summable, so the logarithms of the factors
are summable, and the product is the exponential of their sum.
\end{proof}

\begin{corollary}[\lean{twinPrime\_of\_ratios}]
If Lee's theorem holds for exponents $(\theta,\nu)$ and the ratios hypothesis
holds for the same exponents, then there are infinitely many twin primes.
\end{corollary}

With Theorem~\ref{thm:c2} the corollary has exactly two hypotheses: a published
argument that is not yet formalised, and an open conjecture. As a check on the
transcription of the Euler product we also prove the degenerate evaluation
$\mathcal{Y}(\alpha,\eta;\eta,\eta)=1$ used in Lee's Lemma~2.2
(\lean{ratiosY\_degenerate}).

\paragraph{What the hypothesis is.} The proof in \cite{lee} uses the ratios
formula only through the mean of a product of two logarithmic derivatives
$\zeta'/\zeta$ with shifts whose imaginary parts grow like a power of the
height. Expanded into primes, that mean is a weighted sum of
$\Lambda(m)\Lambda(n)$ over pairs with $m/n$ close to $1$, and the uniformity in
the imaginary parts is what isolates a single difference $h$. Our reading is
that the hypothesis is the prime-pair conjecture with a power saving, written in
the language of zeta, rather than an independent fact about zeros. It is not
known to follow from the Riemann Hypothesis. Under RH, $\zeta'/\zeta$ to the
right of the critical line is a sum over zeros, the mean becomes a double sum
over pairs of zeros, and the term that carries $\mathfrak{S}(h)$ comes entirely
from correlations between \emph{different} zeros \cite{montgomery,bk}. RH fixes
the real parts of the zeros and says nothing about those correlations.

\paragraph{Numerics.} \lean{experiments/ratios\_test.py} evaluates the ratios
formula to height $2\cdot10^4$. For real shifts the remainder stays bounded
(about $0.2$, a relative size of $10^{-5}$). Once the shifts have imaginary
parts of size $1$ to $3$ the remainder behaves like a random walk, with
median increment exponents $0.38$ and $0.41$ and maximum $0.47$ over $80$ random
shift sets, all below the threshold $(39-\sqrt{177})/32\approx0.803$ that \cite{lee} requires.
This supports the hypothesis at this height and shows where its difficulty
lies; it tests bounded imaginary shifts only, not the growing ones the proof
uses. \lean{experiments/twin\_zeros.py} finds the first $100$ zeta zeros in the
twin prime count to $10^9$, with amplitude relative to the prime count of
$2.33+0.79i$ (standard error $0.20$), against $4C_2\approx2.64$ predicted and
$2C_2\approx1.32$ for a random thinning of the primes to the same density. The
shortfall and the phase lag are unexplained.

\section{Discussion}\label{sec:discussion}

Three probes end at one place.

\begin{itemize}
\item \emph{Zeta zeros.} The input needed is a correlation between zeros, not
their location.
\item \emph{Residue decoys.} The input needed is the parity of $\Omega(p+2)$.
\item \emph{Bilinear sums.} The input needed is a parity sum over $p+2$.
\end{itemize}

Each is a correlation of $\Lambda$ or $\lambda$ with a shift of itself: the
Hardy--Littlewood and Chowla family. None of this is new as mathematics; the
parity phenomenon has been understood since Selberg. What we add is (a) a
machine-checked exact reformulation whose generality over arbitrary sets makes
it a test, (b) explicit sets at a concrete size that pass every local test and
have no pairs, with code, and (c) a formal statement of the cleanest conditional
route we know, with its elementary parts checked.

A proof would need a family of integers in which both $p$ and $p+2$ lie inside
some structure that makes a parity-type bilinear sum provable. We know of no
such structure and we do not propose one.

\paragraph{What the Riemann Hypothesis archive contains.} The announcement
\cite{rhclaim} describes its argument only in outline. As far as we can follow
it, the target is Weil's positivity criterion, and the method descends from
Stepanov's proof for function fields, which builds a function vanishing to high
multiplicity. The integers have no derivative with which to count multiplicity,
and the announcement says the substitute is a chain of theta lifts between
unitary groups, each contributing one unit of multiplicity, of length about
$\sqrt{X}$ for a sum of $\Lambda(n)$ over $n\le X$. We do not understand the
proof, and its author says he does not yet understand all of it. On 21 November
2026 we downloaded the archive ($175{,}190$ files) and searched the text of its
Lean sources, without building it. The five headline theorems concern only the
zeros of $\zeta$ and of Dirichlet $L$-functions. We found no occurrence of twin
primes, of the Hardy--Littlewood conjecture, of the Liouville function or
$\Omega$, and no product of $\Lambda$ or $\mu$ with a shift of itself. Every
occurrence of the name Chowla is the Chowla--Selberg formula. The M\"obius
function appears in mollifiers, zero-detecting sums and M\"obius inversion. A
text search cannot exclude a statement written in notation we did not think of;
the patterns are listed in \lean{lean/IDEAS.md}, idea~8.

One question seems worth recording, and it is a question, not a claim. Over
$\mathbb{F}_q[T]$ the twin prime conjecture is a theorem for suitable $q$: Sawin
and Shusterman \cite{sawinshusterman} show that on certain subspaces the
M\"obius function is mimicked by Dirichlet characters, which removes the parity
obstruction there. As we understand their mechanism, it rests on
differentiation in characteristic $p$, where the derivative of a $p$-th power
vanishes. That is the same missing tool, a derivative on the integers, that the
announcement claims to have replaced. Whether the replacement, if it is sound,
could carry any part of their argument to $\mathbb{Z}$ we cannot judge.

For anyone, human or machine, who believes they have an argument: express it as
a function that, given a finite set of integers in $(x,\tfrac{51}{50}x)$,
returns a lower bound for the number of pairs at distance $2$, and run it on the
set $E$ of Section~\ref{sec:parity}. If the bound is positive, the argument is
wrong.

\section*{Reproducibility}

Repository: \url{https://github.com/oceanseth/primefold}. The Lean project is in
\lean{lean/} (Lean \lean{v4.33.1}, Mathlib revision \lean{0df444a3}); build with
\lean{lake exe cache get \&\& lake build}. The experiments are single Python
files in \lean{experiments/} with their outputs stored next to them as
\lean{results-*.txt}. All numbers in this note are copied from those files.

\section*{Statement on AI use}

The Lean proofs, the experiments and the first draft of this text were written
by the language model Claude Fable~5.2 working under the direction of the
author, who is responsible for the content. The complete record of prompts and
responses is published in the repository as \lean{llm-turn-history.jsonl}. The
formal statements can be checked independently of any trust in the model by
building the Lean project; the numerical claims can be checked by running the
scripts. The informal reading of \cite{lee} in Section~\ref{sec:chain} has no
such check and should be treated accordingly.

\begin{thebibliography}{99}

\bibitem{bk} E.~B. Bogomolny and J.~P. Keating, Gutzwiller's trace formula and
spectral statistics: beyond the diagonal approximation, \emph{Phys. Rev. Lett.}
\textbf{77} (1996), 1472--1475.

\bibitem{chowla} S. Chowla, \emph{The Riemann Hypothesis and Hilbert's Tenth
Problem}, Gordon and Breach, 1965.

\bibitem{conreysnaith} J.~B. Conrey and N.~C. Snaith, Applications of the
$L$-functions ratios conjectures, \emph{Proc. London Math. Soc.} (3) \textbf{94}
(2007), 594--646.

\bibitem{fi} J. Friedlander and H. Iwaniec, The polynomial $X^2+Y^4$ captures
its primes, \emph{Ann. of Math.} (2) \textbf{148} (1998), 945--1040.

\bibitem{hl} G.~H. Hardy and J.~E. Littlewood, Some problems of `Partitio
Numerorum'; III: On the expression of a number as a sum of primes, \emph{Acta
Math.} \textbf{44} (1923), 1--70.

\bibitem{lee} S. Lee, The Hardy--Littlewood asymptotic for prime pairs of fixed
even difference from a uniform zeta-ratios conjecture, preprint,
arXiv:2610.32672v2, 2026.

\bibitem{mr} K. Matom\"aki and M. Radziwi\l\l, Multiplicative functions in short
intervals, \emph{Ann. of Math.} (2) \textbf{183} (2016), 1015--1056.

\bibitem{mathlib} The mathlib Community, The Lean mathematical library, in
\emph{Proc. 9th ACM SIGPLAN Int. Conf. on Certified Programs and Proofs
(CPP 2020)}, 367--381.

\bibitem{maynard} J. Maynard, Small gaps between primes, \emph{Ann. of Math.}
(2) \textbf{181} (2015), 383--413.

\bibitem{montgomery} H.~L. Montgomery, The pair correlation of zeros of the zeta
function, in \emph{Analytic Number Theory}, Proc. Sympos. Pure Math. \textbf{24},
Amer. Math. Soc., 1973, 181--193.

\bibitem{polymath} D.~H.~J. Polymath, Variants of the Selberg sieve, and bounded
intervals containing many primes, \emph{Res. Math. Sci.} \textbf{1} (2014),
Art.~12.

\bibitem{rhclaim} L. Alp\"oge, Lean archive claiming a formal proof of the
Riemann Hypothesis, announced 20 November 2026 at
\url{https://alpo.ge/remark.pdf}; no written proof at the time of writing.
Third-party mirror and notes:
\url{https://github.com/Deicyde/riemann-hypothesis}.

\bibitem{rhcheck} A.~V. Sutherland, independent build and kernel re-checks of
\cite{rhclaim}, \url{https://math.mit.edu/~drew/grh/}, as updated 21 November
2026.

\bibitem{sawinshusterman} W. Sawin and M. Shusterman, On the Chowla and twin
primes conjectures over $\mathbb{F}_q[T]$, preprint, arXiv:1808.04001.

\bibitem{selberg} A. Selberg, On elementary methods in primenumber-theory and
their limitations, in \emph{Den 11te Skandinaviske Matematikerkongress
(Trondheim, 1949)}, 13--22.

\bibitem{tao} T. Tao, The logarithmically averaged Chowla and Elliott
conjectures for two-point correlations, \emph{Forum Math. Pi} \textbf{4} (2016),
e8.

\end{thebibliography}

\end{document}
